Following the channel-subset results above \ref{sec:exp1_single_channel}, we further tested whether a single electrode was sufficient as a self-contained detector by running the complete Stage A--Stage C Proposed-Med pipeline separately on each single-electrode selection in both Internal and Cao2018, using 30-second epochs and the median centre. Internal comprised 46 sessions and 6 single-electrode selections (Fp1, Fp2, AF3, AF4, F3 and F4), whereas Cao2018 comprised 58 sessions and 4 single-electrode selections (Fp1, Fp2, F3 and F4). Because each selection contained exactly one channel, these values required no best-channel-per-session oracle, unlike every other condition in the article, and so corresponded directly to what a deployed single-electrode system would produce.

Under single-electrode operation, Proposed-Med retained near-reference $F_1$ only at the frontopolar pair, whereas performance fell sharply at the other electrodes (Table~\ref{tab:exp1_solo_vs_montage}; Figure~\ref{fig:exp1_single_channel}). On Internal, the best single-electrode result was obtained at Fp1 with $F_1=83.73\%$, corresponding to 94.9\% of the full-montage reference, while Fp2 reached $F_1=83.37\%$ and AF4 and AF3 fell to $F_1=75.16\%$ and $F_1=67.59\%$, respectively; the weakest electrode was F3 at $F_1=39.85\%$, corresponding to 45.2\% of the full-montage reference. On Cao2018, Fp1 was also best with $F_1=77.65\%$, corresponding to 96.3\% of the full-montage reference, whereas F3 and F4 reached $F_1=49.30\%$ and $F_1=49.03\%$, respectively, with F4 the weakest at 60.8\% of the full-montage reference. After Bonferroni correction, scoring the same electrode inside the full montage rather than alone did not significantly increase $F_1$ at Internal Fp1 ($\Delta F_1=+2.59$ percentage points, $p=1.000$) or Fp2 ($\Delta F_1=+3.45$ percentage points, $p=1.000$), or at Cao2018 Fp1 ($\Delta F_1=+2.03$ percentage points, $p=1.000$) or Fp2 ($\Delta F_1=+2.86$ percentage points, $p=0.991$). In contrast, the increase was significant at every weaker electrode: Internal AF4 ($\Delta F_1=+5.16$ percentage points, $p=0.001$), AF3 ($\Delta F_1=+6.17$ percentage points, $p=0.002$), F3 ($\Delta F_1=+5.79$ percentage points, $p=0.008$), and F4 ($\Delta F_1=+5.26$ percentage points, $p<0.001$), and Cao2018 F3 ($\Delta F_1=+10.21$ percentage points, $p<0.001$) and F4 ($\Delta F_1=+8.96$ percentage points, $p<0.001$). For these weaker electrodes, the gain was concentrated in recall rather than precision: for Internal AF4, recall increased from 70.10\% alone to 77.67\% in the full montage while precision changed from 87.65\% to 86.91\%, and for Cao2018 F3, recall increased from 43.23\% to 54.95\% while precision changed from 71.32\% to 72.64\%.

Consistent with the electrode-level comparison, collapsing the single-electrode selections to their region mean provided a second, coarser comparison, and because every evaluated single electrode in both corpora belonged to the frontal region, this yielded one frontal-region comparison per corpus rather than four regions (Table~\ref{tab:region_performance}; Table~\ref{tab:exp1_subset_summary}). On Internal, the single-electrode-alone region mean was $F_1=64.95\%$, compared with $F_1=69.68\%$ for the same electrodes scored within the full-montage run, a difference of $\Delta F_1=-4.74$ percentage points; a single session-paired Wilcoxon signed-rank test gave $p<0.001$ and an effect size of $r=-0.80$. On Cao2018, the corresponding values were $F_1=62.89\%$ for the single-electrode-alone region mean and $F_1=68.91\%$ within the full-montage run, a difference of $\Delta F_1=-6.02$ percentage points, $p<0.001$, and $r=-0.72$. Each corpus therefore required only one Wilcoxon signed-rank comparison at this coarse level, so no Bonferroni factor was applied. For comparison, running the same frontal electrodes together as one self-contained multi-electrode subset produced $F_1=66.92\%$ on Internal and $F_1=66.61\%$ on Cao2018, higher than the corresponding single-electrode-alone region means but below the same electrodes' full-montage region means of $F_1=69.68\%$ and $F_1=68.91\%$, respectively.
